\documentclass[letterpaper, 10pt]{article}
\usepackage{graphicx}
\usepackage[hidelinks]{hyperref}
\usepackage[top=0.5in, bottom=0.5in, left=0.7in, right=0.7in]{geometry}
\usepackage{kotex}
\pagestyle{empty}

\begin{document}

%----------HEADING----------
\begin{center}
    \textbf{\LARGE Kim Jihwan} \\ \vspace{2pt}
    \href{mailto:wlghks0333@g.skku.edu}{\underline{wlghks0333@g.skku.edu}} $|$
    \href{https://github.com/Elarge2025}{\underline{github.com/Elarge2025}}
\end{center}

%-----------TECHNICAL SKILLS-----------
\section*{Technical Skills}
\hrule
\vspace{4pt}
\begin{itemize}
    \item[] \textbf{Programming}: C, C++, Python
    \item[] \textbf{Frameworks \& Tools}: Git, GitHub, VS Code
    \item[] \textbf{Languages}: Korean (Native), English (Intermediate)
\end{itemize}

%-----------EXPERIENCE-----------
\section*{Experience}
\hrule
\vspace{4pt}
\begin{itemize}
    \item[]
        \textbf{Member} \hfill Mar. 2025 -- Dec. 2025 \\
        \textit{LikeLion at Sungkyunkwan University} \hfill \textit{Suwon, Korea}
        \begin{itemize}
            \item 웹 개발 기초 및 프로젝트 실습 수료
            \item HTML, CSS, JavaScript 기반 서비스 개발 참여
        \end{itemize}
\end{itemize}

%-----------SELECTED PROJECTS-----------
\section*{Selected Projects}
\hrule
\vspace{4pt}
\begin{itemize}
    \item[]
        \textbf{LikeLion Hackathon} \hfill Summer 2025 \\
        \textit{Sungkyunkwan University}
        \begin{itemize}
            \item 팀 프로젝트로 웹 서비스 개발
            \item 프론트엔드 담당, HTML/CSS/JavaScript를 활용한 UI 구현
        \end{itemize}
\end{itemize}

%-----------HONORS & AWARDS-----------
\section*{Honors \& Awards}
\hrule
\vspace{4pt}
\begin{itemize}
    \item[] \textbf{SQLD} (SQL Developer) \hfill 2025
    \item[] \textbf{ADsP} (Associate Data Analysis Professional) \hfill 2025
\end{itemize}

%-----------EDUCATION-----------
\section*{Education}
\hrule
\vspace{4pt}
\begin{itemize}
    \item[]
        \textbf{Sungkyunkwan University} \hfill Mar. 2024 -- Present \\
        \textit{Suwon, Korea}
        \begin{itemize}
            \item 제1전공: 철학과 (Department of Philosophy)
            \item 복수전공: 소프트웨어학과 (Department of Software)
        \end{itemize}
\end{itemize}

\end{document}